\begin{proof}

\begin{enumerate}[label={$\langle1\rangle$\arabic*.},leftmargin=*]
\item Write $D\subset\widetilde S$ for the conductor support and $N\subset S$ for the singular curve. Normalization is a homeomorphism off these sets, so additivity of the topological Euler characteristic gives $e(S)=e(\widetilde S)-e(D)+e(N)$. In the Hopf case $e(\widetilde S)=0$, $D$ is a disjoint union of two elliptic curves, and $N$ is one elliptic curve; all three Euler numbers vanish. The ruled case has the same conductor curves and $e(\widetilde S)=e(\PP^1)e(E)=0$.

\item In the rational case the normalization is a Hirzebruch surface blown up twice, so its Euler number is $4+2=6$. The conductor is a cycle of six rational curves, with six nodes, so $e(D)=12-6=6$. Its image consists of three rational curves with exactly two triple points; counting the three preimages of each triple point once gives $e(N)=6-2(3-1)=2$. Hence $e(S)=2$.

\item Away from finitely many critical values the proper holomorphic map is a differentiable fibre bundle, and a two-torus has Euler number zero. Additivity for a stratification by this open base and the exceptional points therefore gives $e(Y)=\sum_c e((Y_c)_{\rm red})$. Nilpotents do not change the underlying topological space. The assumed fibre classes and \eqref{cdp:eq:A.18} give \eqref{cdp:eq:A.19}. The equality with $c_3$ is the complex-oriented Gauss--Bonnet theorem. Poincare duality for the connected oriented six-manifold gives $e(Y)=2-2b_1+2b_2-b_3$. Under the two Betti vanishings this is $2-b_3$, so $b_3=2-2r\ge0$ and $r$ is zero or one.

\end{enumerate}
\end{proof}
